\documentclass[UTF8]{ctexart}
\usepackage{geometry}
\usepackage[backend=biber, style=authoryear]{biblatex} % 示例配置
\addbibresource{references.bib} % 替换为你的参考文献文件
\geometry{margin=1.5cm, vmargin={0pt,1cm}}
\setlength{\topmargin}{-1cm}
\setlength{\paperheight}{29.7cm}
\setlength{\textheight}{25.3cm}

% useful packages.
\usepackage{amsfonts}
\usepackage{amsmath}
\usepackage{amssymb}
\usepackage{amsthm}
\usepackage{enumerate}
\usepackage{graphicx}
\usepackage{multicol}
\usepackage{fancyhdr}
\usepackage{layout}
\usepackage{listings}
\usepackage{xcolor}
\lstset{language=C++}
\setlength{\headheight}{13pt}  % 增加头部高度
\usepackage{float, caption}

% Listings setting for code formatting
\lstset{
    basicstyle=\ttfamily,
    keywordstyle=\color{blue},
    commentstyle=\color{gray},
    stringstyle=\color{red},
    breaklines=true,
    frame=single,
    numbers=left,
    numberstyle=\tiny,
    stepnumber=1,
    tabsize=4
}

% some common commands
\newcommand{\dif}{\mathrm{d}}
\newcommand{\avg}[1]{\left\langle #1 \right\rangle}
\newcommand{\difFrac}[2]{\frac{\dif #1}{\dif #2}}
\newcommand{\pdfFrac}[2]{\frac{\partial #1}{\partial #2}}
\newcommand{\OFL}{\mathrm{OFL}}
\newcommand{\UFL}{\mathrm{UFL}}
\newcommand{\fl}{\mathrm{fl}}
\newcommand{\op}{\odot}
\newcommand{\Eabs}{E_{\mathrm{abs}}}
\newcommand{\Erel}{E_{\mathrm{rel}}}

\begin{document}

\pagestyle{fancy}
\fancyhead{}
\lhead{陈梓锐}
\chead{22435036}
\rhead{\today}


%\begin{abstract}
 %   The abstract is not necessary for the theoretical homework, 
   % but for the programming project, 
    %you are encouraged to write one.      
%\end{abstract}

% ============================================
\subsection*{如何使用}
在code中build文件下使用make命令编译，编译成功后，在build文件夹下会生成可执行文件P,在test目录下的test.josn调整初始条件即可使用不同的测试用例。

2维情形下运行较慢，因为在插值时对向量进行了多次复制，导致效率较低。

\subsection*{Problem $exp(y+sin(x))$}
\subsubsection*{1维VCycle}
规则边界，Dirichlet边界条件，fullWeighting松弛, linear插值。








% ===============================================

\end{document}